\documentclass[11pt,a4paper]{article}

\usepackage[T1]{fontenc}
\usepackage[utf8]{inputenc}
\usepackage[magyar]{babel}
\usepackage{lmodern}
\usepackage[margin=2.5cm]{geometry}
\usepackage{amsmath,amssymb}
\usepackage{bm}
\usepackage{graphicx}
\usepackage{booktabs}
\usepackage{xcolor}
\usepackage{enumitem}
\usepackage[hidelinks]{hyperref}

\graphicspath{{figures/}}

% Jelölő a még hiányzó részekhez
\newcommand{\todo}[1]{\textcolor{red}{\textbf{[TODO: #1]}}}
\newcommand{\E}{\mathbb{E}}
\newcommand{\Prob}{\mathbb{P}}
\newcommand{\indep}{\perp\!\!\!\perp}

% A szimulációs kód a GitHubon (commithez kötött link, így a sorszámok nem csúsznak el)
\newcommand{\kodurl}{https://github.com/KMateJ/psm-beadando/blob/1777c530520dd1b15733ab8f1e7eae4fac5e2dac/psm_szimulacio.Rmd?plain=1}
\newcommand{\kodlink}[2]{\href{\kodurl\##1}{\textcolor{blue}{#2}}}

% Forgatókönyvenként egy ábrasor (figures/abra-X-1, az Rmd generálja), alatta a kódrészek linkjei
% #1: kód, #2: cím, #3: az adatgeneráló folyamat sorai, #4: a becsült modellek sorai
\newcommand{\abra}[4]{%
  \begin{figure}[htbp]
    \centering
    \includegraphics[width=\textwidth]{abra-#1-1.png}
    \caption{#2. Balra: becsült PS egy mintában; középen: |SMD| párosítás előtt és után
      (szaggatott vonal: 0,1); jobbra: a becslések eloszlása az 5000 ismétlésben (szaggatott vonal: $\tau = 2$).}
    {\footnotesize Kód: \kodlink{#3}{adatgeneráló folyamat} \textbullet{}
      \kodlink{#4}{becsült modellek} \textbullet{} \kodlink{L76-L90}{párosítás} \textbullet{}
      \kodlink{L277-L358}{ábra}}
  \end{figure}}

\setlength{\parindent}{0pt}
\setlength{\parskip}{0.6em}

\title{\textbf{Propensity score matching:}\\\textbf{módszer, alkalmazások és korlátok}}
\author{%
  {Nagy Levenete} ({DU38XD}, {levente.nagy1@stud.uni-corvinus.hu})\\
  {Kocsis Máté József} ({P6FT0Y}, {mate.kocsis@uni-corvinus.hu})\\
  {Juhász Ferenc} ({LCRXUR}, {ferenc.juhasz@stud.uni-corvinus.hu})\\
  {Kolozsvári Bence Tamás} ({CIQV92}, {bence.kolozsvari@stud.uni-corvinus.hu})
  }
  
\date{2026. október 13.}

\begin{document}
\maketitle

%======================================================================
\section*{1. Bevezetés}

A közgazdaságtan egyik alapvető célja különböző beavatkozások (\emph{treatment}) okozati hatásának,
az átlagos kezelési hatásnak (\emph{Average Treatment Effect}, ATE) a mérése. A standard megoldás
a randomizált kontrollált kísérlet (\emph{randomized controlled trial}, RCT), ahol a kezelés kiosztása
véletlenszerű, így a kezelt és a kontrollcsoport várhatóan minden megfigyelt és megfigyelhetetlen
jellemzőjében megegyezik. A gyakorlatban azonban sokszor csak megfigyeléses adatok állnak
rendelkezésre, ahol a kezelés felvétele nem véletlen, hanem a résztvevők saját döntésén vagy egyéb
tényezőkön múlik.

Ilyenkor kiválasztási torzítás (\emph{selection bias}) lép fel, a kezelésre és a kimenetelre egyaránt ható
zavaró változók (\emph{confounder}-ek) miatt. Ha egyszerűen összehasonlítjuk a kezelt és a nem kezelt
csoport átlagos kimenetelét, torzított eredményt kapunk, mert a különbség egy része a csoportok
eredendő különbségeiből fakad.
 Bár a regressziós modellek képesek kontrollálni ezeket a változókra, sok változó esetén ez nehézkes ( dimenzió átka,
\emph{curse of dimensionality}). A \emph{propensity score matching} (PSM) ezt a problémát úgy kezeli, hogy a confoundereket egyetlen változóba, a kezelés valószínűségébe sűríti, és e szerint teszi összehasonlíthatóvá a
kezelt és a kontroll megfigyeléseket.

%======================================================================
\section*{2. Módszertan matematikai megalapozása}

Minden egyénre ($i = 1, \dots, N$) adott a $T_i \in \{0, 1\}$ kezelési indikátor, a megfigyelt confounderek
$X_i \in \mathbb{R}^p$ vektora, valamint két potenciális kimenetel, $Y_i(1)$ és $Y_i(0)$. A megfigyelt kimenetel
\[
  Y_i = T_i\, Y_i(1) + (1 - T_i)\, Y_i(0).
\]
A célparaméterek az átlagos kezelési hatás, illetve a kezeltekre vonatkozó átlagos kezelési hatás
(\emph{Average Treatment Effect on the Treated}, ATT):
\[
  \mathrm{ATE} = \E\bigl[Y_i(1) - Y_i(0)\bigr], \qquad
  \mathrm{ATT} = \E\bigl[Y_i(1) - Y_i(0) \mid T_i = 1\bigr].
\]
A megfigyeléses adatokon számolt naiv különbség torzított:
\begin{align*}
  \Delta &= \E[Y_i \mid T_i = 1] - \E[Y_i \mid T_i = 0] \\
         &= \underbrace{\E\bigl[Y_i(1) - Y_i(0) \mid T_i = 1\bigr]}_{\text{ATT}}
          + \underbrace{\E\bigl[Y_i(0) \mid T_i = 1\bigr] - \E\bigl[Y_i(0) \mid T_i = 0\bigr]}_{\text{szelekciós torzítás}}.
\end{align*}
A naiv becslés tehát csak akkor egyezik meg az ATT-vel, ha a szelekciós torzítás nulla. Ez
randomizált kísérletben garantált, megfigyeléses adatokban a confounderek miatt általában nem
teljesül.
A megfigyeléses adatokból oksági hatás becsléséhez két feltevés szükséges:
\begin{enumerate}[nosep]
  \item \textbf{Feltételes függetlenség} (\emph{unconfoundedness}): $\bigl(Y_i(1), Y_i(0)\bigr) \indep T_i \mid X_i$.
  \item \textbf{Közös tartomány} (\emph{common support/overlap}): $0 < \Prob(T_i = 1 \mid X_i = x) < 1$ minden $x$-re.
\end{enumerate}
A propensity score a kezelés valószínűsége a megfigyelt tulajdonságok mellett:
\[
  e(x) = \Prob(T = 1 \mid X = x).
\]
\textbf{Tétel (Rosenbaum–Rubin, 1983)}: ha teljesül a feltételes függetlenség, akkor
\begin{enumerate}[nosep]
  \item (kiegyensúlyozó tulajdonság, \emph{balancing property}) $X \indep T \mid e(X)$,
  \item $\bigl(Y(1), Y(0)\bigr) \indep T \mid e(X)$.
\end{enumerate}
Vagyis elég az $e(X)$ skalárra kondicionálni, nem kell az egész $X$ vektorra.
A propensity score becslése egy osztályozási probléma. Azt  valószínűséget akarjuk megbecsülni, hogy egy adott személy – a megfigyelt tulajdonságai alapján – megkapja-e a kezelést vagy sem.
Ezt leggyakrabban logisztikus regresszióval becsüljük:
\[
  e(X_i) = \Prob(T_i = 1 \mid X_i)
  = \frac{1}{1 + \exp\bigl(-(\beta_0 + \beta_1 X_{1i} + \dots + \beta_k X_{ki})\bigr)} \in (0, 1).
\]
Ezután minden kezelt egyént párosítunk egy hasonló score-ú nem kezelttel. A két legelterjedtebb párosítási
algoritmus a legközelebbi szomszéd (\emph{nearest-neighbour}) és a kernel matching. A párosítás
után ellenőrizni kell, hogy a minta valóban kiegyensúlyozott lett-e, azaz a párosítás eltüntette-e a csoportok közötti különbségeket az alapjellemzőkben. Ha igen, a párosított csoportok
kimeneteleinek különbsége valóban becsli a kezelés okozati hatását.

%======================================================================
\section*{3. Szakirodalmi példák}

\textbf{1. Cikk:} Dehejia, R. H., \& Wahba, S. (1999). \emph{Causal effects in nonexperimental studies:
Reevaluating the evaluation of training programs.}

Azt vizsgálták, hogy egy amerikai államilag támogatott munkaerőpiaci tréning növeli-e a résztvevők
jövőbeli jövedelmét. A probléma eredete, hogy a 70-es évek közepén az USA-ban indítottak egy
drága RCT-t. Sorsolással döntöttek, ki kap képzést, ki boldogul egyedül a munkaerőpiacon,
az RCT kimutatta a valós pozitív hatást. 1986-ban egy közgazdász vette ennek a National
Supported Work program adatait, kidobta a kísérlet valódi kontrollcsoportját, helyette vett egy
országos megfigyelésen alapuló adatbázist. Lineáris regresszióval próbálta kiszámolni a tréning
valódi hatását, kezelt embereket átlag amerikaiakhoz hasonlítva, negatív hatást kapott. Az
ellentmondás feloldása az volt, hogy a linkelt tanulmány szerzőpárosa logisztikus regressziót
épített, hogy megbecsüljék, mekkora esélye (propensity score) lett volna egy átlag amerikainak
arra, hogy bekerüljön a tréningbe, felhasználva a confoundereket (életkor, iskolai végzettség,
korábbi jövedelem\dots). A nearest-neighbour matching során minden tréningezett személyhez hasonló propensity score-ú, a kontrollcsoportba tartozó személyt párosítottak. A korábbi torzítást az okozta, hogy a tréningre főleg az alacsonyabb jövedelmi és társadalmi percentilisekből érkeztek, ezt nem volt figyelembe véve. Miután a propensity score alapján kizárták a nem összehasonlítható megfigyeléseket, a becsült hatás alig tért el eredeti tanulmány kimutatott nagy pozitív hatásától.

\textbf{2. Cikk:} Morgan, S. L. (2001). \emph{Counterfactuals, causal effect heterogeneity, and the Catholic
school effect on learning.} Sociology of Education, 341--374.

Amerikában az volt a közmegegyezés, hogy a katolikus (magán) középiskolák jobb oktatást
biztosítanak. A katolikus iskolába járó gyerekek kompetencia tesztjei magasan verték az állami
iskolásokét. A politikusok ebből azt a következtetést vonták le, hogy az állami iskolák sokkal
rosszabbak a magániskoláknál. Értelemszerűen ezt nem lehetett RCT-vel tesztelni. Egy gyerek
nem véletlenül kerül magániskolába. A propensity score kiszámításához használt tényezők
(socioeconomic status, középiskola előtti teszteredmények stb.) együttesen meghatározták, egy
gyerek mekkora eséllyel köt ki egy katolikus iskolában. A szerző a logisztikus regresszióval kapott
propensity score alapján nem átlagot vont (ATE), hanem sorba állította a diákokat PS alapján,
és csoportokra bontotta őket. Ez a stratification matching. A high propensity csoport (tehetős
szülők, egyetemi elvárások) egyéneire a magániskola hozzáadott értéke elhanyagolható volt. A
low propensity csoport (szegény család, gyenge korábbi jegyek) tagjainak kis esélyük van fizetős
iskolába jutni, azonban bejutás esetén a hatás kiemelkedően pozitív volt. Megmutatkozott a hatás
heterogenitása.

\textbf{3. Cikk:} Gilligan, M. J., \& Sergenti, E. J. (2008). \emph{Do UN interventions cause peace? Using
matching to improve causal inference.}

Az ENSZ békefenntartás kapcsán a naiv statisztikák azt mutatták, hogy ahol az ENSZ jelen van, ott
hamarabb kiújul a háború. Ebből politikusok azt szűrték le, hogy a kéksisakosok inkompetensek. A
szerzők vettek 117 polgárháborút, és logisztikus regresszióval kiszámolták, mekkora valószínűséggel
menne oda az ENSZ (confounderek: konfliktus hossza, ország megosztottsága, háború előtti
halálesetek száma). A propensity score matching után, mikor az ENSZ missziókat kizárólag
magukra hagyott, de ugyanolyan véres háborúkkal hasonlították össze, kiderült, hogy az ENSZ
beavatkozással a béke átlagosan sokkal tovább tart. A common support szerepe: azokat a konfliktusokat, amelyekben az ENSZ egyáltalán nem avatkozik be (jellemzően a könnyen kezelhető esetek), nincs hozzájuk hasonló kezelt eset, ezért kizárja őket az elemzésből. AZ ENSZ-missziók nem véletlenül történnek, jellemzően nehezebb konfliktusokba küldik őket, ezért volt itt csúnya félrevezetés.

%======================================================================
\section*{4. Szimuláció és eredmények}
Minden forgatókönyvben a confounderek független standard normális eloszlásúak, a kezelés logit modellből
származik, és a kezelés valódi hatása konstans, $\tau = 2$ (így ATE = ATT = 2):
\[
  T_i \sim \text{Bernoulli}\bigl(\Lambda(\text{index}_i)\bigr), \qquad
  Y_i = y_0(X_i) + \tau T_i + \varepsilon_i, \quad \varepsilon_i \sim N(0,1).
\]

ahol $\Lambda(\cdot)$ a standard logisztikus függvény, $y_0(X_i)$ a bázis kimenetel, az $\text{index}_i$ pedig a confounderek lineáris (vagy nemlineáris) kombinációja.

A propensity score-t ($e(X_i)$) logisztikus regresszióval becsüljük. Ezután a PS logitján ($L_i = \text{logit}(e(X_i))$) 1:1 legközelebbi szomszéd párosítást (nearest-neighbor matching) végzünk visszatevéssel, mindkét irányban.
Minden kezeltnek megkeressük a legközelebbi kontrollt: 
\[
j(i) = \arg\min_{j: T_j = 0} |L_i - L_j|
\]
Minden kontrollnak megkeressük a legközelebbi kezeltet: 
\[
i(j) = \arg\min_{i: T_i = 1} |L_j - L_i|
\]
Az így kapott egyéni hatások átlaga az ATE-t becsli
\cite{abadie2006}. Viszonyítási alap a naiv átlagkülönbség. A kiegyensúlyozottságot a standardizált
átlagkülönbséggel (SMD) mérjük, a szokásos küszöb $|\mathrm{SMD}| < 0{,}1$. Forgatókönyvenként 5000
Monte Carlo ismétlést futtatunk $n = 10\,000$ megfigyeléssel. Az eredményeket az 1.~táblázat, a
diagnosztikai ábrákat az 1--5.~ábra mutatja.

\subsection*{4.1 Első adatgeneráló folyamat: a PSM működik}

\textbf{A) Jól specifikált PS-modell.} $\text{index} = 0{,}8X_1 - 0{,}6X_2 + 0{,}5X_3$ és
$y_0 = 1{,}5X_1 - X_2 + X_3$. Mindhárom confounder megfigyelt, és a becsült logit modell megegyezik a
valódival. A naiv különbség 1,83-mal felülbecsüli a hatást, mert a kezeltek $X_1$ és $X_3$ szerint
kedvezőbb helyzetűek. A párosítás után minden |SMD| 0,1 alatt van, a PSM átlagos becslése 2,004,
a szórása 0,028. A propensity score itt a confounderek teljes torzító hatását kiszűri.

\subsection*{4.2 Második adatgeneráló folyamat: a PSM nem működik}

\textbf{B) Rosszul specifikált PS-modell.} A kezelés és a kimenetel is négyzetesen függ a confounderektől:
$\text{index} = 0{,}6(X_1^2 - 1) + 0{,}6(X_2^2 - 1) + 0{,}3X_3$, $y_0 = 1{,}2X_1^2 + 1{,}2X_2^2 + X_3$.
A szokásos lineáris logit ezt nem tudja leírni: a becsült PS alig jelzi előre a kezelést
(AUC = 0,57), és a becslés torzítása 1,96, alig kevesebb, mint a naiv különbségé (2,17). Az $X_1, X_2, X_3$
változók a párosítás után kiegyensúlyozottnak tűnnek, csak a modellből kimaradt $X_1^2$ és $X_2^2$
mutatja a problémát. Ha a négyzetes tagokat is bevonjuk, a torzítás 0,03-ra csökken. A hiba tehát
javítható, de csak akkor, ha felismerjük a rossz modellválasztást.

\textbf{C) Meg nem figyelt confounder.} $\text{index} = 0{,}6X_1 - 0{,}5X_2 + U$ és
$y_0 = X_1 - X_2 + 1{,}5U$, ahol $U$ nem megfigyelt, így sérül a feltételes függetlenség.
A párosítás után $X_1$ és $X_2$ tökéletesen kiegyensúlyozott, az $U$ viszont nem
(3.~ábra), a torzítás 1,25 marad. Ha $U$ megfigyelhető lenne, a torzítás 0,005 volna. Ez a legtrükkösebb eset a tesztek látszólagos egyensúlyt mutatnak, és a probléma az adatokból
nem mutatható ki.

\textbf{D) Gyenge overlap.} Az A) eset négyszer erősebb szelekcióval, így a PS-eket 0 és az 1 felé tolja (AUC = 0,96). Kevés a hasonló pár, ugyanazokat a megfigyeléseket sokszor használjuk fel: a
torzítás 0,175, a szórás pedig hatszorosa az A) esetének (0,175 vs.\ 0,028). A PS logitján 0,2 szórásnyi calipert alkalmazunk
\cite{austin2011} a torzítást 0,01-re csökkenti, de csak a megfigyelések 87\%-át tartja meg. 

\textbf{E) Instrumentum a PS-modellben.} $\text{index} = 0{,}6X_1 - 0{,}5X_2 + 3Z$ és
$y_0 = X_1 - X_2$, tehát $Z$ csak a kezelésre hat. A $Z$ bevonása javítja az előrejelzést
(AUC = 0,92), torzítást nem okoz (2,01), de rontja az overlapet, így a becslés szórása hatszorosára nő
(0,145 vs.\ 0,024). A lényeg az, hogy a jó PS-modell célja nem a kezelés tökéletes előrejelzése, hanem a rejtett torzítások kiszűrése. \cite{brookhart2006}.

\begin{table}[htbp]
\centering\small
\caption{Monte Carlo eredmények (5000 ismétlés, $n = 10\,000$, valódi hatás $\tau = 2$)}
\begin{tabular}{llrrrrr}
\toprule
 & Becslő & Átlag & Torzítás & Szórás & RMSE & AUC \\
\midrule
A) jól specifikált    & naiv különbség           & 3,833 & 1,833 & 0,042 & 1,834 &       \\
                      & PSM                      & 2,004 & 0,004 & 0,028 & 0,029 & 0,759 \\
\addlinespace
B) rossz specifikáció & naiv különbség           & 4,172 & 2,172 & 0,052 & 2,172 &       \\
                      & PSM, lineáris logit      & 3,960 & 1,960 & 0,059 & 1,961 & 0,569 \\
                      & PSM, helyes specifikáció & 2,029 & 0,029 & 0,036 & 0,046 & 0,746 \\
\addlinespace
C) rejtett confounder & naiv különbség           & 3,977 & 1,977 & 0,041 & 1,978 &       \\
                      & PSM, $U$ nélkül          & 3,254 & 1,254 & 0,043 & 1,255 & 0,670 \\
                      & PSM, $U$-val             & 2,005 & 0,005 & 0,031 & 0,031 & 0,782 \\
\addlinespace
D) gyenge overlap     & naiv különbség           & 5,050 & 3,050 & 0,034 & 3,051 &       \\
                      & PSM                      & 2,175 & 0,175 & 0,175 & 0,247 & 0,957 \\
                      & PSM, caliper 0,2$^{*}$   & 2,013 & 0,013 & 0,102 & 0,103 & 0,957 \\
\addlinespace
E) instrumentum       & naiv különbség           & 2,493 & 0,493 & 0,035 & 0,495 &       \\
                      & PSM, $Z$-vel             & 2,010 & 0,010 & 0,145 & 0,146 & 0,923 \\
                      & PSM, $Z$ nélkül          & 2,001 & 0,001 & 0,024 & 0,024 & 0,599 \\
\bottomrule
\end{tabular}

\smallskip
\footnotesize $^{*}$ A megfigyelések 87\%-a marad meg a párosítás után.
\end{table}

\abra{A}{A) jól specifikált PS-modell}{L120-L126}{L159-L163}
\abra{B}{B) rosszul specifikált PS-modell}{L128-L136}{L164-L171}
\abra{C}{C) meg nem figyelt confounder}{L138-L144}{L172-L179}
\abra{D}{D) gyenge overlap}{L146-L147}{L180-L187}
\abra{E}{E) instrumentum a PS-modellben}{L149-L155}{L188-L195}

%======================================================================
\section*{5. Következtetések és limitációk}

A szimulációk alapján a PSM akkor ad torzítatlan becslést, ha (1) minden confounder megfigyelt
(nincs rejtett torzítás), (2) a PS-modell helyesen specifikált, és (3) a kezelt és a kontrollcsoport
PS-eloszlása kellően átfed. Ha ezek teljesülnek, a PSM a naiv összehasonlítás teljes szelekciós
torzítását eltünteti (A eset).

AHa viszont ezek a feltételek sérülnek, a következmények eltérőek lehetnek. A rossz specifikáció javítható (rugalmasabb modell,
magasabb rendű tagok és interakciók), de ezt észre kell venni. A kiegyensúlyozottsági
ellenőrzés csak a modellbe bevont változókra mond valamit, ezért érdemes a négyzetes tagokra és az
interakciókra is elvégezni (B eset). A meg nem figyelt confounder okozta torzítást viszont az adatokból sem
kimutatni, sem kiszűrni nem lehet (C eset). Ez a PSM legnagyobb korlátja. Gyenge overlap esetén a becslés zajos és torzított. A
caliper ezen segít, de ilyenkor elveszítjük a minta egy részét (D eset). Végül a
PS-modellbe a confoundereket kell bevonni, nem a kezelés legjobb előrejelzőit: a csak a kezelésre ható
változók torzítást nem okoznak, de jelentősen növelik a becslés szórását (E eset). Az eredmények
a párosítási paraméterektől (visszatevés, caliper, párok száma, párosítási algoritmus) is függenek.

Szimulációnk egyszerűsítő feltevései: konstans kezelési hatás, normális eloszlású confounderek és logit
kezelési modell. Heterogén hatásnál az ATE és az ATT eltér, és a caliper vagy a common support miatti
kizárás a becsült paramétert is megváltoztatja. Lehetséges általánosítások: regresszió helyett rugalmasabb, gépi tanulási algoritmusokra bízzuk a modellépítést, párosítás helyett inverz valószínűségi súlyozást (IPW) használunk.

%======================================================================
\begin{thebibliography}{9}\small

\bibitem{abadie2006}
Abadie, A., \& Imbens, G. W. (2006). Large sample properties of matching estimators for average treatment
effects. \emph{Econometrica}, 74(1), 235--267.

\bibitem{austin2011}
Austin, P. C. (2011). Optimal caliper widths for propensity-score matching when estimating differences in
means and differences in proportions in observational studies. \emph{Pharmaceutical Statistics}, 10(2), 150--161.

\bibitem{brookhart2006}
Brookhart, M. A., Schneeweiss, S., Rothman, K. J., Glynn, R. J., Avorn, J., \& Stürmer, T. (2006). Variable
selection for propensity score models. \emph{American Journal of Epidemiology}, 163(12), 1149--1156.

\bibitem{dehejia1999}
Dehejia, R. H., \& Wahba, S. (1999). Causal effects in nonexperimental studies: Reevaluating the
evaluation of training programs. \emph{Journal of the American Statistical Association}, 94(448), 1053--1062.

\bibitem{gilligan2008}
Gilligan, M. J., \& Sergenti, E. J. (2008). Do UN interventions cause peace? Using matching to improve
causal inference. \emph{Quarterly Journal of Political Science}, 3(2), 89--122.

\bibitem{lalonde1986}
LaLonde, R. J. (1986). Evaluating the econometric evaluations of training programs with experimental
data. \emph{American Economic Review}, 76(4), 604--620.

\bibitem{morgan2001}
Morgan, S. L. (2001). Counterfactuals, causal effect heterogeneity, and the Catholic school effect on
learning. \emph{Sociology of Education}, 74(4), 341--374.

\bibitem{rosenbaum1983}
Rosenbaum, P. R., \& Rubin, D. B. (1983). The central role of the propensity score in observational
studies for causal effects. \emph{Biometrika}, 70(1), 41--55.

\end{thebibliography}

\end{document}
